\documentclass[10pt,a4paper]{amsart}
\usepackage[margin=19mm]{geometry}
\usepackage{amsmath,amssymb,mathtools}
\usepackage[scr=rsfs,cal=euler]{mathalfa}
\usepackage{xfrac}
\usepackage[unicode,hidelinks,bookmarks=false]{hyperref}
\newcommand{\ie}{i.e.\ }
\newcommand{\cf}{cf.\ }
\newcommand{\resp}{respectively\ }
\newcommand{\Subsection}{\subsection}
\providecommand{\nobreakdash}{\nobreak\hbox{-}}
\setlength{\emergencystretch}{2em}
\multlinegap0pt
\usepackage{bm}
\usepackage{bbm}
\usepackage{dsfont}
\usepackage{xspace}
\usepackage{cancel}
\usepackage{mathtools}
\usepackage{amsthm}
\makeatletter
\@ifundefined{prop}{\newtheorem{prop}{Proposition}}{}
\@ifundefined{thm}{\newtheorem{thm}{Theorem}}{}
\makeatother
\title[Generalized capillary-rise models: existence and fast solver]{Generalized capillary-rise models: existence and fast solvers in integral H\"older spaces}
\author{Josefa Caballero, {\L}ukasz P{\l}ociniczak, Kishin Sadarangani}
\date{Source: arXiv:2510.05801v1 (CC BY 4.0)}
\begin{document}
\maketitle

\noindent\emph{Source and licence.} Generalized capillary-rise models: existence and fast solvers in integral H\"older spaces. Version arXiv:2510.05801v1 (CC BY 4.0), distributed under the Creative Commons Attribution 4.0 International licence, as recorded in the arXiv licence metadata for this exact version and shown on the abstract page https://arxiv.org/abs/2510.05801v1. Full rights notice: licenses/arxiv-2510.05801-CC-BY-4.0.txt.\par\smallskip

\noindent\textit{Editorial scope.} Counted unit: the Thm whose source label is \texttt{thm:LinearCollocationConvergence} in the article (source0.tex lines 583--589, source label \texttt{thm:LinearCollocationConvergence}), delivered together with the article's own complete original proof of it (source0.tex lines 590--657, one \texttt{proof} environment of 68 lines). No mathematics is written, repaired or re-derived: statement and proof are the article's own text line for line. Required context retained uncounted: eqn:MainEq (lines 113--115); eqn:AssumG (lines 117--120); eqn:AssumF (lines 124--127); eqn:OpeartorT (lines 168--170); prop:TNorm (lines 242--251); eqn:AssumRho (lines 244--246); eqn:AssumR0 (lines 273--275); thm:Existence (lines 325--327); eqn:LinearInterpolation (lines 392--394); prop:LInftyJInterpolation (lines 477--486); eqn:LinearCollocationScheme (lines 515--521). Every definition, display and label of those lines is retained unchanged. Omitted from this package and not redistributed: 1--112, 116--116, 121--123, 128--167, 171--241, 247--272, 276--324, 328--391, 395--476, 487--514, 522--582, 658--1157. Nothing inside a retained excerpt is omitted or altered: every retained body is delivered exactly as the article's lines are written, so no omission receipt of any kind accompanies this package. Citations: the retained excerpts carry no citation command at all, so no bibliography entry of the article is transcribed and no reference is redistributed. Reference adaptation: none was needed and none was made; every reference inside the delivered unit resolves inside it, so no unresolved reference or citation is produced. Printed slips kept literal: no printed word, symbol, hypothesis or number of the counted statement or of its proof was altered. Layout and class: the article's own class is \texttt{article} with packages including amssymb, amsmath, amsthm, algorithm, algpseudocode, fullpage, babel, graphicx, booktabs, xcolor, enumitem, hyperref, autonum; the wrapper is the lane's pinned \texttt{amsart} editorial wrapper at 10pt on A4, which restates the theorem-like environments the retained text needs and adds the article's own macro definitions that the retained text needs (none), with extra wrapper packages (bm, bbm, dsfont, xspace, cancel, mathtools). The delivered numbering is the unit's own. Assets: no retained span contains a figure environment, a graphics-inclusion command, a PDF-page inclusion, a table or a TikZ picture; no image file is copied into this package, the package declares no asset dependency and no image file is redistributed with it. Licence: version arXiv:2510.05801v1 of this article is distributed under the Creative Commons Attribution 4.0 International licence, as recorded in the arXiv licence metadata for this exact version and shown on the abstract page https://arxiv.org/abs/2510.05801v1. A full rights notice is delivered at licenses/arxiv-2510.05801-CC-BY-4.0.txt. Characters: every retained excerpt is pure ASCII, so the delivered engine, XeLaTeX with its default Latin Modern text font, reports no missing character.
\begin{equation}\label{eqn:MainEq}
	x(t) = \int_0^t G(t,s) f(s, x(s)) ds, \quad t \in[0,1].
\end{equation}

\begin{equation}\label{eqn:AssumG}
	\tag{H1}
	G \in C([0,1]^2), \quad M_G := \max_{s,t\in[0,1]} |G(s,t)|, \quad G(\cdot, s) \in C_\rho[0,1], \quad \|G(\cdot, s)\|_\rho \leq K_G, \quad s \in [0,1],
\end{equation}

\begin{equation}\label{eqn:AssumF}
	\tag{H2}
	|f(t,x) - f(t, y)| \leq \phi(|x-y|) \quad \text{for any} \quad t \in [0,1] \quad \text{and}\quad M_f := \sup\left\{|f(t,0)|: \, t\in[0,1]\right\} < \infty,
\end{equation}

\begin{equation}\label{eqn:OpeartorT}
	(Tx)(t) := \int_0^t G(t,s) f(s, x(s)) ds, \quad t \in [0,1].
\end{equation}

\begin{prop}\label{prop:TNorm}
	Suppose that $0<\alpha\leq 1$ and $\alpha \leq \beta < \infty$. Assume that
	\begin{equation}\label{eqn:AssumRho}
		\int_0^1 \sigma^{-(\beta+1)} \rho(\sigma)^\frac{\beta}{\alpha} d\sigma < \infty.
	\end{equation}
	Then, for any $x \in J_{\alpha,\beta}[0,1]$ we have $Tx \in J_{\alpha,\beta}[0,1]$ and
	\begin{equation}\label{eqn:TNorm}
		\|Tx\|_{\alpha,\beta} \leq \left(\phi(\|x\|_{\alpha,\beta}) + M_f\right)  \left(\int_0^1 \sigma^{-(\beta+1)} \left(K_G\rho(\sigma) + \sigma M_G \right)^\frac{\beta}{\alpha} d\sigma\right)^\frac{\alpha}{\beta}.
	\end{equation}
\end{prop}

	\begin{equation}\label{eqn:AssumR0}
		\left(\phi(r_0) + M_f\right)  \left(\int_0^1 \sigma^{-(\beta+1)} \left(K_G\rho(\sigma) + \sigma M_G \right)^\frac{\beta}{\alpha} d\sigma\right)^\frac{\alpha}{\beta} < r_0.
	\end{equation}

\begin{thm}\label{thm:Existence}
	Assume \eqref{eqn:AssumRho} and suppose that there exists $r_0>0$ satisfying \eqref{eqn:AssumR0}. Then, the integral equation \eqref{eqn:MainEq} has a solution in $J_{\alpha,\beta}[0,1]$ with $0<\alpha<1$ and $\alpha\leq \beta<\infty$. 
\end{thm}

\begin{equation}\label{eqn:LinearInterpolation}
	I_h x(t) = \frac{1}{h} \left(x(t_{n+1}) (t- t_n) + x(t_n) (t_{n+1}- t)\right), \quad t \in [t_{n}, t_{n+1}], \quad 0\leq n \leq N-1, \quad x \in C[0,1].
\end{equation}

\begin{prop}\label{prop:LInftyJInterpolation}
	Let $x \in J_{\alpha,\beta}[0,1]$. Then, for its linear interpolation \eqref{eqn:LinearInterpolation} and $0<\alpha\leq \beta$ we have
	\begin{equation}
		\|x - I_h x\|_\infty \leq \left(\frac{\beta}{2^{\beta}-1}\right)^\frac{\alpha}{\beta} \left(1-2^{-\frac{\beta}{\frac{\beta}{\alpha}-1}}\right)^{-\frac{\beta - \alpha}{\beta}}h^\alpha \left(\int_0^h \sigma^{-\beta-1} \omega(x,\sigma)^\frac{\beta}{\alpha}\right)^\frac{\alpha}{\beta},
	\end{equation}
	while for $\beta\rightarrow\infty$ we obtain
	\begin{equation}
		\|x-I_h x\|_\infty \leq \frac{h^\alpha}{2^{\alpha}-1} \sup_{0<\sigma\leq h}\frac{\omega(x,\sigma)}{\sigma^\alpha}.
	\end{equation}
\end{prop}

\begin{equation}\label{eqn:LinearCollocationScheme}
	\begin{cases}
		x_h \text{ is a linear function on each } [t_n, t_{n+1}], & 0 \leq n < N, \\
		x_h(t_n^+) = x_h(t_{n+1}^-), & 0 < n < N, \\
		x_h(t_n) = (Tx_h)(t_n), & 0 \leq n \leq N, \\
	\end{cases}
\end{equation}

\begin{thm}[Convergence of the linear collocation scheme]\label{thm:LinearCollocationConvergence}
	Let the assumptions of Theorem \ref{thm:Existence} be satisfied. Suppose that $\phi$ from \eqref{eqn:AssumF} satisfies $\phi(x) \leq L x$ with $L>0$. Then, we have
	\begin{equation}\label{eqn:LinearCollocationError}
		\|x - x_h\|_\infty \leq C_{G,f} \left(\frac{\beta}{2^{\beta}-1}\right)^\frac{\alpha}{\beta} \left(1-2^{-\frac{\beta}{\frac{\beta}{\alpha}-1}}\right)^{-\frac{\beta - \alpha}{\beta}}h^\alpha \left(\int_0^h \sigma^{-\beta-1} \omega(x,\sigma)^\frac{\beta}{\alpha}\right)^\frac{\alpha}{\beta},
	\end{equation}
	for some constant $C_{G,f}>0$ dependent only on $G$ and $f$. 
\end{thm}

\begin{proof}
	By the definition of the collocation scheme \eqref{eqn:LinearCollocationScheme} and the interpolation operator \eqref{eqn:LinearInterpolation}, since $x_h$ is a piecewise linear polynomial, we have for any $t\in[0,1]$
	\begin{equation}
		x_h(t) = I_h x_h(t) = \sum_{n=1}^N L_n(t) x_h(t_n) = \sum_{n=1}^N L_n(t) Tx_h(t_n) = I_h Tx_h (t).
	\end{equation}	
	Let $e_h := x - x_h$ be the error of the scheme. Then, by the above, we have
	\begin{equation}
		e_h(t) = x(t) - x_h(t) = Tx(t) - I_h Tx_h(t) = Te_h(t) + Tx_h(t) - I_h T x_h(t).
	\end{equation}
	Therefore, by the definition of the operator $T$ as in \eqref{eqn:OpeartorT} and the regularity assumptions \eqref{eqn:AssumG}-\eqref{eqn:AssumG}, we have
	\begin{equation}
		|e_h(t)| \leq K_G \int_0^t \phi(|e_h(s)|) ds + \rho_h(t),
	\end{equation}
	where 
	\begin{equation}
		\rho_h(t) := |Tx_h(t) - I_h T x_h(t)|. 
	\end{equation}
	But, due to the assumption $\phi(|e_h(s)|) \leq L |e_h(s)|$, and hence
	\begin{equation}
		|e_h(t)| \leq K_G L \int_0^t |e_h(s)|ds + \rho_h(t).
	\end{equation} 
	Now, by the classical Gronwall inequality, for some new constant $D_{G,f}$ we obtain
	\begin{equation}
		|e_h(t)| \leq \rho_h(t) + D_{G,f} \int_0^t \rho_h(s) ds \leq (1 + D_{G,f}) \|\rho_h\|_\infty. 
	\end{equation}
	Hence, setting $C_{G,f} := 1+ D_{G,f}$,
	\begin{equation}
		\|e_h\|_\infty \leq C_{G,f} \|\rho_h\|_\infty. 
	\end{equation}	
	But due to Proposition \ref{prop:TNorm} we have $Tx \in J_{\alpha,\beta}[0,1]$, and therefore from Proposition \ref{prop:LInftyJInterpolation} we can obtain the estimate on the interpolation error $\|\rho_h\|_\infty$. The proof concludes. 
	%Use the linear interpolation operator $I_h$ and decompose the error as follows 
	%\begin{equation}
	%	\|x-x_h\|_\infty \leq \|x - I_h x\|_\infty + \|I_h x - x_h\|_\infty =: \|\rho_h\|_\infty + \|\theta_h\|_\infty.
	%\end{equation}
	%By Theorem \ref{thm:Existence} we know that $x\in J_{\alpha,\beta}[0,1]$ hence, the estimate on the first term $\rho_h$ follows from Proposition \ref{prop:LInftyJInterpolation} 
	%\begin{equation}\label{eqn:RhoBound}
	%	\rho_h \leq \left(\frac{\beta}{2^{\beta}-1}\right)^\frac{\alpha}{\beta} \left(1-2^{-\frac{\beta}{\frac{\beta}{\alpha}-1}}\right)^{-\frac{\beta - \alpha}{\beta}}h^\alpha \left(\int_0^h \sigma^{-\beta-1} \omega(x,\sigma)^\frac{\beta}{\alpha}\right)^\frac{\alpha}{\beta},
	%\end{equation}
	%Therefore, we only have to focus on the second term $\theta_g$ in which we compare two piecewise linear functions with the same number of degrees of freedom. In the Lagrange basis, we can write
	%\begin{equation}
	%	I_h x(t) = \sum_{n=1}^N x(t_n) L_n(t), \quad x_h(t) = \sum_{n=1}^N x_h(t_n) L_n(t).
	%\end{equation}
	%therefore,
	%\begin{equation}
	%	|I_h x(t) - x_h(t)| = |\theta_h(t)| \leq \sum_{n=1}^N L_n(t) \left|x(t_n) - x_h(t_n)\right|.
	%\end{equation}
	%But, by the main equation \eqref{eqn:MainEq} and the collocation condition \eqref{eqn:LinearCollocationScheme} we have
	%\begin{equation}
	%\begin{split}
	%	|\theta_h(t)| 
	%	&\leq \sum_{n=1}^N |L_n(t)| \int_0^t |G(t, s)| \left|f(s, x(s)) - f(s, x_h(s))\right|ds \\
	%	&\leq K_G \sum_{n=1}^N L_n(t) \int_0^t \left|\phi(x(s) - x_h(s))\right|ds \leq K_G \int_0^t \left|\phi(x(s) - x_h(s))\right|ds,
	%\end{split}
	%\end{equation}
	%where we have used \eqref{eqn:LinearInterpolationConstant} the regularity assumptions on $G$ and $f$ as in \eqref{eqn:AssumG} and \eqref{eqn:AssumG}. Now, by subadditivity of $\phi$ and since $t \leq 1$, we further obtain
	%\begin{equation}
	%	|\theta_h(t)| \leq K_G L \|\rho_h(s)\|_\infty + K_G\int_0^t |\theta_h(s)| ds.
	%\end{equation}
	%Invoking the classical Gronwall's lemma (see []) further yields
	%\begin{equation}
	%	|\theta_h(t)| \leq C_{G,f} \|\rho_h(s)\|_\infty,
	%\end{equation}
	%with some constant $C_{G,f} > 0$. Therefore, by taking supremum over $t\in [0,1]$ and going back to the error decomposition, we obtain
	%\begin{equation}
	%	\|e_h\|_\infty \leq (1+C_{G,f}) \|\rho_h\|_\infty. 
	%\end{equation}
	%Combining this with the bound for $\rho_n$ as in \eqref{eqn:RhoBound} concludes the proof. 
\end{proof}

\end{document}
